% !TEX root = ../../proposal.tex

\section{Conclusion}
\label{sec:conclusion}
Phase 1 produces four concrete outputs for the broader research programme:

\begin{enumerate}

    \item \textbf{A curated concept-pair benchmark} of 24 pairs spanning the four theory-grounded taxonomy groups, with shared-noise evaluation protocols and per-pair metadata (taxonomy group, theoretical basis, and expected failure mode).
    This benchmark provides a reusable testbed for evaluating compositional generation methods beyond this research.

    \item \textbf{A trajectory measurement framework}: shared-noise MDS projections and per-step MSE metrics that characterise the gap geometrically and temporally.
    The framework is model-agnostic and can be applied to any pair of composition conditions.

    \item \textbf{A temporal characterisation} of gap onset and amplification: the decomposition of cumulative divergence into early-step (layout) and late-step (identity resolution) contributions.
    This identifies the denoising interval where guidance schemes would need to intervene to close the gap.

    \item \textbf{A reachability map}: per-category cycle-consistency scores under SD-IPC inversion, providing lower-bound evidence for the structural inaccessibility of logical composition endpoints.

\end{enumerate}

Together, these outputs constitute the principled measurement framework called for by Research Goal~1 and provide the empirical grounding for Hypothesis~1.

Phase~2 builds on this with three concrete outputs:

\begin{enumerate}

    \item \textbf{A four-regime controlled design:} standard dSprites serves as a positive
    control, correlated dSprites provides the failure setting, Regime~C contributes both benchmark
    disentanglement models and matched-supervision controls, and correlated dSprites with SP-VAE
    provides the structural intervention. This design isolates the effect of latent factorisation
    while keeping the composition test fixed.

    \item \textbf{A held-out composition-success evaluation:} attribute classifiers trained on
    disjoint data verify whether a composed sample recovers both target factors from withheld
    orientation--scale combinations, supported by representative qualitative examples.

    \item \textbf{A benchmarked representation comparison:} when the data distribution induces
    entanglement, composition degrades; benchmark disentanglement models test how much of that loss
    can be recovered through independence objectives alone, matched-supervision controls test how
    much is recovered from supervision without latent partitioning, and SP-VAE tests whether
    explicit factor-addressable partitioning improves further. This comparison is the core
    empirical support for H2.

\end{enumerate}

Phase~3 translates these findings to the clinical domain:

\begin{enumerate}

    \item \textbf{Clinical held-out composition evaluation:} pathology presence classifiers, a
    co-morbidity fidelity classifier trained on real co-morbid CXR cases, and spatial plausibility
    checks constitute the clinical analogue of the Phase~2 evaluation logic.

    \item \textbf{An empirical answer to H3:} whether approximate anatomical independence between
    cardiomegaly and pleural thickening is sufficient for factorised composition to produce
    ecologically valid co-morbid chest X-rays not present in the training data.

\end{enumerate}

The unifying claim across all three phases is the same: compositional generation of
out-of-distribution combinations works when the internal representations of the composed concepts
do not conflict with each other. Phase~1 documents that this conflict exists and is predictable
from the structure of the training distribution. Phase~2 shows that this conflict can be reduced
in a controlled setting by improving latent factorisation under a fixed composition rule. Phase~3
tests whether the approximate anatomical separation of disease conditions in chest X-rays provides
the same leverage in a real clinical setting --- and thereby determines whether factorised
composition can produce clinically valid co-morbid presentations without ever having seen them in
training.

